% ---------------------------------------------------------------------------
%  QML Lab research project 03: Projected vs Fidelity Kernels
%  Build:  latexmk -pdf main.tex   (needs qmlproject.sty, harvardqml.bst, references.bib)
% ---------------------------------------------------------------------------
\documentclass[11pt,a4paper]{article}
\usepackage{qmlproject}
\qmlheader{Project 03: Projected vs Fidelity Kernels}

\hypersetup{pdftitle={QML Lab project 03: Projected vs Fidelity Kernels}}

\begin{document}

\begin{center}
  {\small\textcolor{black!55}{QML Lab \textbullet\ research project 03}}\\[0.5em]
  {\LARGE\bfseries Projected vs Fidelity Kernels}\\[0.6em]
  {\large How fast do quantum kernels concentrate as qubits are added, and can projected kernels delay it?}\\[0.6em]
  {\small Living document \textbullet\ last updated \today}
\end{center}

\begin{projectcard}
\small
\begin{tabularx}{\linewidth}{@{}>{\bfseries}l X >{\bfseries}l X@{}}
Quadrant  & \CQ\ (classical data, quantum processor) & Track & Kernels \\
Difficulty & Hard & Tools & PennyLane, NumPy \\
Builds on & Modules 6 and 8 of the QML Lab & Workspace & \path{qml_projects/projected-vs-fidelity-kernels/} \\
\end{tabularx}
\end{projectcard}

\paragraph{Summary.}
Quantum kernels are attractive because they use exponentially large feature spaces, but that same
size makes kernel values concentrate: all pairs of points start to look equally (dis)similar. This
project measures concentration for fidelity kernels and projected kernels from 2 to 12 qubits, for
different feature maps and bandwidths, and with finite measurement shots. The result is a map of
where quantum kernels remain usable.

% ===========================================================================
\section{Background and key concepts}
% ===========================================================================
\subsection{Fidelity and projected kernels}
The fidelity kernel compares whole states, $\kappa^{F}(x,x')=\Tr[\rho(x)\rho(x')]$. The projected
kernel of \citet{huang2021power} compares only reduced single-qubit states $\rho_k(x)$:
\begin{equation}
  \kappa^{PQ}(x,x')=\exp\!\Big(-\gamma\sum_{k=1}^{n}\lVert\rho_k(x)-\rho_k(x')\rVert_F^2\Big),
  \qquad
  \lVert\rho_k-\rho_k'\rVert_F^2=\tfrac12\sum_{P\in\{X,Y,Z\}}\big(\langle P_k\rangle-\langle P_k\rangle'\big)^2 .
\end{equation}
It needs only single-qubit expectation values, which are cheap to measure, and it behaves like an
RBF kernel on a learned, quantum-generated feature vector.

\subsection{Exponential concentration}
\citet{thanasilp2024} prove that for expressive embeddings, entangling embeddings, global
measurements or noise, the kernel values concentrate around a fixed value with variance
\begin{equation}
  \operatorname{Var}_{x,x'}\big[\kappa(x,x')\big]\in O(b^{-n}),\qquad b>1 .
\end{equation}
With $M$ shots, a kernel entry is estimated with standard error about $\sqrt{\kappa(1-\kappa)/M}$, so
once the spread of kernel values falls below this error the Gram matrix is dominated by shot noise.
Projected kernels avoid the global overlap but also concentrate when reduced states approach the
maximally mixed state.

\subsection{Bandwidth and inductive bias}
Scaling features by a bandwidth $\lambda$ keeps encoded states closer together and delays
concentration \citep{shaydulin2022}. \citet{kubler2021} argue more generally that useful quantum
kernels must restrict themselves to a small, task-relevant part of Hilbert space. A related warning
is that settings that provably avoid concentration are often the ones that classical computers can
simulate \citep{cerezo2023simulability}.

% ===========================================================================
\section{Research questions}
% ===========================================================================
The project is open-ended: these questions are the starting point, not the finish line. Add new
ones to the README and the journal as they come up.

\begin{rqlist}
  \item At which qubit count does the off-diagonal variance of each kernel drop below the shot-noise level for $M=10^3$ and $10^4$ shots?
  \item Does reducing the bandwidth delay concentration equally for fidelity and projected kernels?
  \item How do the results depend on the feature map (IQP/ZZ versus a deep hardware-efficient embedding)?
\end{rqlist}

% ===========================================================================
\section{Method}
% ===========================================================================
\subsection{Data}
Synthetic inputs drawn uniformly from $[0,2\pi]^n$ to isolate the effect of the embedding, and a
real data set (breast cancer, PCA to $n$ features) to see the effect on classification. For each
qubit number $n=2,\dots,12$, 200 random pairs estimate the kernel statistics.

\subsection{Models}
Exact state-vector simulation of the IQP/ZZ map and of a hardware-efficient embedding with $L=n$
layers. Fidelity kernels from state overlaps; projected kernels from single-qubit Pauli expectations
with $\gamma$ set by the median heuristic. Finite-shot versions sample from the exact probabilities.
For classification, both kernels feed a precomputed-kernel SVM.

\subsection{Baselines}
An RBF kernel on the same inputs (with tuned width) for the classification part; the analytic
$O(2^{-n})$ reference curve for Haar-random states.

\subsection{Metrics}
\begin{itemize}
  \item Mean and variance of off-diagonal kernel values versus $n$ (log scale).
  \item Effective rank of the Gram matrix.
  \item Test accuracy of the SVM and geometric difference to the RBF kernel.
\end{itemize}

% ===========================================================================
\section{Planned experiments}
% ===========================================================================
Each experiment becomes one \texttt{[[experiments]]} entry in the workspace, so its runs, parameters
and metrics are stored and compared in the portal.

\begin{table}[h]
\centering\small
\begin{tabularx}{\linewidth}{@{}l>{\raggedright}X>{\raggedright}X>{\raggedright\arraybackslash}X@{}}
\toprule
ID & Experiment & What varies & Main metrics \\
\midrule
  E1 & Concentration curves & $n=2..12$, both kernels, IQP map & off-diagonal variance \\
  E2 & Bandwidth & $\lambda$ sweep at each $n$ & variance, accuracy \\
  E3 & Feature map & IQP vs hardware-efficient & variance slope \\
  E4 & Finite shots & $M\in\{10^2,10^3,10^4\}$ & rank, accuracy \\
  E5 & Classification & breast cancer, $n=2..10$ & accuracy, $g$ \\
\bottomrule
\end{tabularx}
\caption{Experiment plan. Results go into the journal and the results log (\cref{sec:log}).}
\end{table}

% ===========================================================================
\section{Running it in the QML Lab}
% ===========================================================================
The workspace lives in \path{qml_projects/projected-vs-fidelity-kernels/}. Scripts import three helpers from
\path{qml_lab/plots.py}:
\begin{itemize}
  \item \texttt{params()} reads the parameters from the portal's run form (defaults when run by hand);
  \item \texttt{save()} stores figures with the run;
  \item \texttt{log\_metrics()} records numbers that appear in the run table.
\end{itemize}

\begin{lstlisting}
qml_projects/projected-vs-fidelity-kernels/
  project.toml          metadata + experiments
  README.md             goal, open questions, ideas
  experiments/          runnable scripts
  docs/main.tex         this document
\end{lstlisting}

First experiment: This experiment is \emph{planned}: write the script, add the block below to \texttt{project.toml}, then run \texttt{python manage.py sync\_projects}.

\begin{lstlisting}[language={}]
[[experiments]]
slug = "concentration"
title = "Kernel concentration from 2 to 12 qubits"
script = "experiments/e01_concentration.py"
figure = "concentration.png"

[experiments.params]
max_qubits = 12
n_pairs = 200
bandwidth = 1.0
feature_map = "iqp"
shots = 0
seed = 0
\end{lstlisting}

\begin{lstlisting}[language=bash]
# from the portal folder: pick up changes to project.toml
python manage.py sync_projects
cd qml_projects/projected-vs-fidelity-kernels
# run by hand with the default parameters
PYTHONPATH=../../qml_lab python experiments/e01_concentration.py
\end{lstlisting}

% ===========================================================================
\section{Quality bar}
% ===========================================================================
\begin{itemize}
  \item A classical baseline with a similar number of trainable parameters, tuned with the same care \citep{bowles2024}.
  \item At least five random seeds per configuration, reported as mean $\pm$ standard deviation.
  \item How results change with the number of qubits.
  \item At least one run under a realistic noise model.
  \item Plot variances on a log scale and fit the slope; report the fitted base $b$ with its uncertainty.
\end{itemize}

% ===========================================================================
\section{Open directions}
% ===========================================================================
\begin{itemize}
  \item $k$-body projected kernels that use reduced states of $k>1$ qubits.
  \item Trainable bandwidth (connects to Project 02).
  \item Check whether the non-concentrating regimes can be simulated classically \citep{cerezo2023simulability}.
\end{itemize}

% ===========================================================================
\section{Risks and limitations}
% ===========================================================================
\begin{itemize}
  \item Conclusions depend on the input distribution; uniform synthetic data is a worst case.
  \item The ``shot-noise threshold'' depends on the chosen $M$; report results for several values.
  \item State-vector memory limits the study to about 14 qubits on a laptop.
\end{itemize}

% ===========================================================================
\section{Results log}
\label{sec:log}
% ===========================================================================
Summarise findings here as they become stable; day-to-day notes belong in the portal journal.

\begin{table}[h]
\centering\small
\begin{tabularx}{\linewidth}{@{}l l >{\raggedright}X >{\raggedright\arraybackslash}X@{}}
\toprule
Date & Run & Finding & Next step \\
\midrule
 & & & \\
 & & & \\
 & & & \\
\bottomrule
\end{tabularx}
\end{table}

\bibliographystyle{harvardqml}
\bibliography{references}

\end{document}
